% !TeX program = pdflatex
% Compilación recomendada (por la bibliografía):
%   pdflatex 798_Lectura_4_2S2026.tex
%   biber documento_apoyo_airflow3
%   pdflatex 798_Lectura_4_2S2026.tex
%   pdflatex 798_Lectura_4_2S2026.tex

\begin{filecontents*}{referencias_airflow3.bib}
@online{airflow_docs,
  title        = {Documentation | Apache Airflow},
  organization = {Apache Software Foundation},
  url          = {https://airflow.apache.org/docs/},
  urldate      = {2026-02-17}
}

@online{airflow_release_notes_3,
  title        = {Release Notes — Airflow Documentation (Airflow 3.0.0 Highlights)},
  organization = {Apache Software Foundation},
  url          = {https://airflow.apache.org/docs/apache-airflow/stable/release_notes.html},
  urldate      = {2026-02-17}
}

@online{airflow_upgrading_3_stable,
  title        = {Upgrading to Airflow 3 — Airflow Documentation},
  organization = {Apache Software Foundation},
  url          = {https://airflow.apache.org/docs/apache-airflow/stable/installation/upgrading_to_airflow3.html},
  urldate      = {2026-02-17}
}

@online{airflow_public_interface,
  title        = {Public Interface for Airflow 3.0+ — Airflow Documentation},
  organization = {Apache Software Foundation},
  url          = {https://airflow.apache.org/docs/apache-airflow/stable/public-airflow-interface.html},
  urldate      = {2026-02-17}
}

@online{airflow_blog_3_ga,
  title        = {Apache Airflow® 3 is Generally Available!},
  organization = {Apache Software Foundation},
  url          = {https://airflow.apache.org/blog/airflow-three-point-oh-is-here/},
  urldate      = {2026-02-17}
}

@online{airflow_upgrade_3_breaking,
  title        = {Upgrading to Airflow 3 (Airflow 3.0.0) — Breaking Changes},
  organization = {Apache Software Foundation},
  url          = {https://airflow.apache.org/docs/apache-airflow/3.0.0/installation/upgrading_to_airflow3.html},
  urldate      = {2026-02-17}
}

@online{airflow_standard_provider_bash,
  title        = {BashOperator — apache-airflow-providers-standard Documentation},
  organization = {Apache Software Foundation},
  url          = {https://airflow.apache.org/docs/apache-airflow-providers-standard/stable/operators/bash.html},
  urldate      = {2026-02-17}
}

@online{airflow_standard_provider_python_api,
  title        = {airflow.providers.standard.operators.python — API Reference},
  organization = {Apache Software Foundation},
  url          = {https://airflow.apache.org/docs/apache-airflow-providers-standard/stable/_api/airflow/providers/standard/operators/python/index.html},
  urldate      = {2026-02-17}
}
\end{filecontents*}

\documentclass[12pt]{article}

% Idioma y codificación
\usepackage[spanish]{babel}
\usepackage[utf8]{inputenc}
\usepackage[T1]{fontenc}
\usepackage{lmodern}

% Página y espaciado
\usepackage{geometry}
\geometry{letterpaper, margin=1in}
\usepackage{setspace}
\onehalfspacing

% Tipografía y microajustes
\usepackage{microtype}
\usepackage{csquotes}

% Figuras (opcional) y tablas
\usepackage{graphicx}
\usepackage{booktabs}
\usepackage{longtable}
\usepackage{array}

% Listas
\usepackage{enumitem}

% Código
\usepackage{listingsutf8}
\usepackage{listings}
\lstset{inputencoding=utf8}
\lstset{literate=%
  {á}{{\'a}}1 {é}{{\'e}}1 {í}{{\'\i}}1 {ó}{{\'o}}1 {ú}{{\'u}}1
  {Á}{{\'A}}1 {É}{{\'E}}1 {Í}{{\'I}}1 {Ó}{{\'O}}1 {Ú}{{\'U}}1
  {ñ}{{\~n}}1 {Ñ}{{\~N}}1 {ü}{{\"u}}1 {Ü}{{\"U}}1
}
\usepackage{xcolor}
\lstdefinestyle{python}{%
  language=Python,
  inputencoding=utf8,
  basicstyle=\ttfamily\small,
  numbers=left,
  numberstyle=\tiny,
  stepnumber=1,
  numbersep=8pt,
  breaklines=true,
  showstringspaces=false,
  tabsize=2,
  frame=single
}
\lstdefinestyle{bash}{%
  language=bash,
  inputencoding=utf8,
  basicstyle=\ttfamily\small,
  numbers=left,
  numberstyle=\tiny,
  stepnumber=1,
  numbersep=8pt,
  breaklines=true,
  showstringspaces=false,
  tabsize=2,
  frame=single
}

% Hipervínculos
\usepackage{hyperref}
\hypersetup{
  colorlinks=true,
  linkcolor=blue,
  urlcolor=blue,
  citecolor=blue
}

% Formato académico
\usepackage{titlesec}
\usepackage{etoolbox}
\titlespacing*{\section}{0pt}{3ex plus 1ex minus .2ex}{2ex}
\titlespacing*{\subsection}{0pt}{2.5ex plus 1ex minus .2ex}{1.5ex}
\titlespacing*{\subsubsection}{0pt}{2ex plus 1ex minus .2ex}{1ex}
\preto\section{\clearpage}

% Encabezados y pies
\usepackage{fancyhdr}
\pagestyle{fancy}
\fancyhf{}
\lhead{Documento de apoyo (Airflow 3.x)}
\rhead{Orquestación de pipelines}
\cfoot{\thepage}
\renewcommand{\headrulewidth}{0.4pt}
\renewcommand{\footrulewidth}{0pt}

% Bibliografía (biblatex + biber)
\usepackage[
  backend=biber,
  style=ieee,
  sorting=nyt,
  maxbibnames=99
]{biblatex}
\addbibresource{../../referencias_airflow3.bib}

\title{\textbf{Documento de apoyo (Sesión 4, Unidad 2)}\\
\large Orquestación de pipelines de datos con Apache Airflow 3.x\\
\normalsize Teoría práctica para entender y usar Airflow en flujos ETL/ELT}
\author{Javier Andrés Valdez González / Auxiliar \\
\small Seminario de Sistemas 2 — Sección P \\
\small Universidad de San Carlos de Guatemala}
\date{Febrero de 2026}

\begin{document}
\maketitle

\begin{abstract}
Este documento complementa el material de la Sesión 3 (ETL y su rol en OLAP) enfocándose en el \textit{cómo operar} un flujo de datos de forma confiable: \textbf{orquestación}. Apache Airflow 3.x es un orquestador de workflows que permite describir procesos como un grafo acíclico dirigido (DAG), ejecutar tareas con control de dependencias, reintentos, observabilidad y gobierno operativo. El énfasis aquí no es redefinir ETL, sino explicar los conceptos, componentes y patrones necesarios para \textbf{autoría, ejecución y operación} de DAGs en Airflow 3.x, considerando los cambios arquitectónicos y de interfaz pública introducidos en la versión 3.
\end{abstract}

\section{Introducción: de ``transformar datos'' a ``operar procesos''}
En un curso de inteligencia de negocios, la construcción de un pipeline rara vez falla por la idea conceptual (extraer--transformar--cargar), sino por su operación diaria: dependencias entre pasos, ejecuciones parciales, reintentos, \textit{backfills}, credenciales, auditoría, trazabilidad, alertas y recuperación ante errores. Airflow se sitúa en ese plano operativo: define \textbf{qué} se ejecuta, \textbf{cuándo} se ejecuta y \textbf{bajo qué condiciones} se considera exitoso.

Airflow 3.0 marca un punto de inflexión al adoptar una arquitectura \textit{service-oriented} y una \textbf{interfaz pública estable} para autores de DAGs basada en \texttt{airflow.sdk} \cite{airflow_release_notes_3,airflow_public_interface}. Por tanto, aprender Airflow 3.x implica dominar: (i) el modelo mental DAG--Task--Run, (ii) la planificación (scheduling) y control de concurrencia, y (iii) prácticas de autoría que sean robustas frente a cambios de versión.

\section{Conceptos esenciales}
\subsection{DAG, tarea y ejecución}
\begin{itemize}[leftmargin=*]
  \item \textbf{DAG (Directed Acyclic Graph):} definición del workflow. Describe tareas y dependencias; no ejecuta por sí mismo.
  \item \textbf{Task:} unidad mínima de trabajo dentro del DAG. En Airflow, una \textit{task} se materializa mediante un \textbf{Operator} (plantilla ejecutable) o mediante la \textbf{TaskFlow API} (\texttt{@task}).
  \item \textbf{DagRun:} una instancia de ejecución del DAG (por horario, manual o evento).
  \item \textbf{TaskInstance:} ejecución específica de una tarea dentro de un \textit{DagRun}. Es lo que realmente ``corre'' y produce logs/estado.
\end{itemize}
En producción, lo importante es separar \textit{definición} (DAG) de \textit{instanciación} (DagRun) y de \textit{ejecución} (TaskInstance). Esta separación explica por qué un DAG puede existir sin ejecutarse, y por qué un mismo DAG genera muchas ejecuciones históricas.

\subsection{Operators, Sensors, Hooks y Providers}
Airflow integra con sistemas externos a través de \textbf{providers} (paquetes versionados de forma independiente del core). Un provider aporta:
\begin{itemize}[leftmargin=*]
  \item \textbf{Operators:} realizan acciones (ejecutar SQL, disparar un job, invocar una API).
  \item \textbf{Sensors:} esperan condiciones (archivo disponible, partición creada). En versiones modernas, muchos sensores pueden ser \textit{deferrable} (no consumen un worker durante la espera).
  \item \textbf{Hooks:} encapsulan conectividad (credenciales, clientes, sesiones) hacia el sistema externo.
\end{itemize}
En Airflow 3, varios operadores comunes se movieron a \texttt{apache-airflow-providers-standard} (por ejemplo \texttt{BashOperator} y \texttt{PythonOperator}), por lo que su importación típica se realiza desde \texttt{airflow.providers.standard...} \cite{airflow_upgrade_3_breaking,airflow_standard_provider_bash}.

\subsection{Comunicación y configuración: XComs, Variables y Connections}
\textbf{XCom} es el mecanismo de comunicación entre tareas para pasar valores pequeños (metadatos, rutas, IDs). No está pensado para transportar dataframes o datasets completos.
\textbf{Variable} y \textbf{Connection} representan configuración y credenciales administradas por Airflow. En Airflow 3 se recomienda acceder a estos recursos mediante el \textbf{Task SDK} (\texttt{airflow.sdk}) o mediante el contexto de tarea \cite{airflow_public_interface}.

\section{Arquitectura de Airflow 3.x (qué corre y por qué importa)}
Airflow 3.x introduce un cambio clave: \textbf{las tareas ya no deben acceder directamente a la base de metadatos}. En su lugar, el \textbf{API server} actúa como punto de acceso central para tareas y workers \cite{airflow_upgrading_3_stable}. Esto reduce riesgos de seguridad (código de usuario con acceso total a metadatos) y mejora escalabilidad (menos conexiones directas al DB).

\subsection{Componentes principales}
\begin{longtable}{p{0.22\textwidth}p{0.74\textwidth}}
\toprule
\textbf{Componente} & \textbf{Responsabilidad operativa}\\
\midrule
\endhead
\textbf{Scheduler} & Decide qué \textit{TaskInstances} están listas (dependencias satisfechas), aplica reglas de concurrencia, encola ejecución y coordina \textit{backfills}. Airflow 3 incluye mejoras relevantes en \textit{backfills} administrados por scheduler \cite{airflow_blog_3_ga,airflow_release_notes_3}.\\
\textbf{DAG Processor} & Proceso dedicado a parsear DAGs. Separa la ``lectura de código'' del resto del sistema, mejorando estabilidad; implica que el \textit{código a nivel de módulo} debe ser liviano. \cite{airflow_upgrade_3_breaking}\\
\textbf{API Server (UI/REST)} & Expone la UI modernizada y la API REST. En Airflow 3, el comando \texttt{airflow api-server} reemplaza al webserver tradicional \cite{airflow_upgrade_3_breaking}.\\
\textbf{Executor / Workers} & Ejecutan las tareas. Airflow soporta distintos ejecutores (Local, Celery, Kubernetes, Edge) y en Airflow 3 se refuerza una \textit{Task Execution API} que permite mayor aislamiento y ejecución remota \cite{airflow_release_notes_3}.\\
\textbf{Metadata DB} & Persistencia de estado (DAGs, runs, tasks, logs referenciados, conexiones/variables en algunos despliegues). En Airflow 3 se busca minimizar acceso directo desde tareas \cite{airflow_upgrading_3_stable}.\\
\bottomrule
\end{longtable}

\subsection{Implicaciones prácticas para autores de DAGs}
\begin{enumerate}[leftmargin=*]
  \item \textbf{No asumir acceso a modelos internos} (SQLAlchemy, sesiones, tablas de Airflow). Está restringido \cite{airflow_upgrading_3_stable}.
  \item \textbf{Usar la interfaz pública estable:} \texttt{airflow.sdk} para DAGs/tareas y utilidades públicas \cite{airflow_public_interface}.
  \item \textbf{Minimizar trabajo en tiempo de parseo:} evitar consultas, lectura de archivos grandes o llamadas de red en el nivel superior del archivo DAG.
\end{enumerate}

\section{Autoría de DAGs en Airflow 3: interfaz estable con \texttt{airflow.sdk}}
\subsection{Imports recomendados}
Airflow 3 formaliza \texttt{airflow.sdk} como interfaz primaria para autores de DAGs: \texttt{dag()}, \texttt{task()}, \texttt{TaskGroup}, \texttt{Asset}, \texttt{Variable}, \texttt{Connection}, y acceso al contexto con \texttt{get\_current\_context()} \cite{airflow_public_interface}.

\begin{lstlisting}[style=python,caption={Estructura mínima de un DAG con TaskFlow API (Airflow 3.x)}]
from datetime import datetime
from airflow.sdk import dag, task

@dag(
    dag_id="ejemplo_minimo_taskflow",
    start_date=datetime(2025, 1, 1),
    schedule="0 2 * * *",
    catchup=False,
    tags=["demo", "airflow3"]
)
def pipeline():
    @task
    def paso_a():
        return "ok"

    @task
    def paso_b(msg: str):
        print(f"Mensaje: {msg}")

    paso_b(paso_a())

pipeline()
\end{lstlisting}

\subsection{Acceso a Variables y Connections (sin DB directo)}
Airflow 3 recomienda acceder a variables y conexiones vía \textbf{contexto} o directamente desde el SDK \cite{airflow_public_interface}.

\begin{lstlisting}[style=python,caption={Acceso a Variable/Connection usando el contexto de tarea (recomendado)}]
from airflow.sdk import task, get_current_context

@task
def usa_configuracion():
    ctx = get_current_context()
    conn = ctx["conn"].get("mi_conn_id")   # Connection administrada por Airflow
    umbral = int(ctx["var"].value.get("umbral_calidad", "0"))  # Variable

    # Ejemplo: usar credenciales para conectar a un servicio externo
    print(conn.host, conn.login, umbral)
\end{lstlisting}

\begin{lstlisting}[style=python,caption={Acceso directo (SDK) a Connection/Variable}]
from airflow.sdk import Connection, Variable

mi_conn = Connection.get("mi_conn_id")
umbral = Variable.get("umbral_calidad", default_var="0")
\end{lstlisting}

\subsection{Operadores del ``standard provider''}
En Airflow 3, para operadores comunes se requiere instalar \texttt{apache-airflow-providers-standard} y usar sus rutas de importación \cite{airflow_upgrade_3_breaking,airflow_standard_provider_bash,airflow_standard_provider_python_api}.

\begin{lstlisting}[style=python,caption={Ejemplo de uso de BashOperator desde el provider estándar}]
from datetime import datetime
from airflow.sdk import dag
from airflow.providers.standard.operators.bash import BashOperator

@dag(dag_id="ejemplo_bash", start_date=datetime(2025, 1, 1), schedule=None, catchup=False)
def ejemplo():
    BashOperator(
        task_id="listar",
        bash_command="ls -la"
    )

ejemplo()
\end{lstlisting}

\section{Planificación: horarios, data intervals, catchup y backfill}
\subsection{Schedule vs. ejecución real}
El atributo \texttt{schedule} define la intención (cron, @daily, o un \textit{timetable}). El scheduler determina cuándo crear \textit{DagRuns}. En Airflow 3, \texttt{catchup} es \textbf{False por defecto}, reduciendo ejecuciones masivas no deseadas al desplegar un DAG nuevo \cite{airflow_upgrade_3_breaking}.

\subsection{Backfills administrados por scheduler}
El \textit{backfill} es la re-ejecución histórica para intervalos pasados (por ejemplo, rehacer particiones faltantes). Airflow 3 mejora el soporte de backfills al administrarlos desde el scheduler, con monitoreo vía UI/API \cite{airflow_blog_3_ga,airflow_release_notes_3}. Operativamente, esto facilita:
\begin{itemize}[leftmargin=*]
  \item control centralizado de concurrencia,
  \item diagnósticos consistentes,
  \item ejecución segura sin scripts manuales externos.
\end{itemize}

\subsection{Versionado de DAG}
Airflow 3 incorpora \textbf{DAG Versioning}: un \textit{DagRun} se completa usando la versión del DAG con la que inició, incluso si se despliega una nueva versión a mitad de ejecución \cite{airflow_blog_3_ga}. Esto reduce resultados inconsistentes por cambios de código durante un run activo, y mejora trazabilidad (código, estructura y logs asociados a una versión).

\section{Control de flujo y confiabilidad (patrones operativos)}
\subsection{Reintentos, timeouts y reglas de disparo}
A nivel de tarea, se recomienda definir parámetros como:
\begin{itemize}[leftmargin=*]
  \item \texttt{retries} y \texttt{retry\_delay}: para fallas transitorias (red, rate limits).
  \item \texttt{execution\_timeout}: evita tareas colgadas.
  \item \texttt{trigger\_rule}: controla cómo se comporta una tarea cuando hay upstream fallando o saltado (\textit{skipped}).
\end{itemize}
Estas configuraciones son parte del diseño de confiabilidad: un pipeline robusto asume fallas parciales y define qué significa ``completo'' de forma explícita.

\subsection{Branching y ShortCircuit}
Para rutas condicionales (por ejemplo, \textit{si no hay nuevos datos, no cargar}), se usa branching (por ejemplo \texttt{BranchPythonOperator}) o \texttt{ShortCircuitOperator}. En Airflow 3 estos operadores están disponibles vía el provider estándar \cite{airflow_standard_provider_python_api}.

\subsection{TaskGroup y legibilidad}
Los DAGs crecen rápido. \texttt{TaskGroup} permite agrupar tareas y mantener legibilidad sin recurrir a SubDAGs (removidos) \cite{airflow_upgrade_3_breaking,airflow_public_interface}.

\subsection{Dynamic Task Mapping (map) como patrón de paralelismo}
Cuando existe una lista de particiones/archivos, es común mapear una tarea sobre esa lista (p.\,ej., un task por archivo). El valor práctico es: paralelismo controlado + observabilidad por partición (cada mapeo es una TaskInstance con \texttt{map\_index}).

\section{Orquestación orientada a datos: Assets y scheduling por eventos}
Airflow 3 refuerza un estilo \textit{data-aware}: en lugar de pensar ``ejecuta a las 2am'', se puede pensar ``ejecuta cuando el dataset X se actualice'' \cite{airflow_blog_3_ga,airflow_release_notes_3}. Este enfoque es útil cuando:
\begin{itemize}[leftmargin=*]
  \item los orígenes publican eventos (colas, almacenamiento, CDC),
  \item el horario fijo no garantiza disponibilidad de datos,
  \item se desea encadenar pipelines por producción/consumo de \textit{assets}.
\end{itemize}

\begin{lstlisting}[style=python,caption={Ejemplo conceptual de asset como contrato de datos (Airflow 3.x)}]
from datetime import datetime
from airflow.sdk import dag, task, asset, Asset

ventas_diarias = Asset("warehouse://dw/fact_ventas_diarias")

@dag(dag_id="producer_ventas", start_date=datetime(2025, 1, 1), schedule="0 2 * * *", catchup=False)
def producer():
    @task(outlets=[ventas_diarias])
    def construir_fact():
        # construir partición y publicarla como asset
        pass
    construir_fact()

@dag(dag_id="consumer_kpis", start_date=datetime(2025, 1, 1), schedule=[ventas_diarias], catchup=False)
def consumer():
    @task
    def recalcular_kpis():
        pass
    recalcular_kpis()

producer()
consumer()
\end{lstlisting}
El ejemplo ilustra el contrato: el productor declara que \textit{produce} un asset (\texttt{outlets}) y el consumidor se agenda cuando el asset está disponible (en lugar de por cron).

\section{Observabilidad: UI, logs y trazabilidad}
Airflow 3 moderniza la UI (React) y mejora la navegación entre vistas \textit{task-centric} y \textit{asset-centric} \cite{airflow_blog_3_ga}. En operación, la observabilidad suele reducirse a tres preguntas:
\begin{enumerate}[leftmargin=*]
  \item \textbf{¿Qué falló y dónde?} (estado de TaskInstances, logs, stacktrace).
  \item \textbf{¿Desde cuándo y con qué impacto?} (runs históricos, backfills, retrasos).
  \item \textbf{¿Qué cambió?} (versionado de DAG, cambios de código, parámetros).
\end{enumerate}
A nivel de diseño, conviene que cada tarea sea \textbf{idempotente} (re-ejecutable), con logs útiles y validaciones explícitas (por ejemplo conteos, checksums o pruebas de calidad).

\section{Despliegue y operación}
Aunque el despliegue suele ser responsabilidad del equipo de plataforma, el autor de DAG debe comprender restricciones básicas:

\subsection{Ejecución local vs. producción}
\begin{itemize}[leftmargin=*]
  \item \textbf{Local:} útil para aprender. Se puede usar el \textit{Docker stack} oficial o un executor local.
  \item \textbf{Producción:} suele requerir múltiples componentes (scheduler, api-server, dag-processor, workers) y un DB robusto. El Helm Chart oficial es una vía común en Kubernetes \cite{airflow_docs}.
\end{itemize}

\subsection{CLI y ``airflowctl''}
Airflow 3 introduce \texttt{airflowctl} como CLI que interactúa exclusivamente con la API REST, evitando acceso directo a la base de metadatos y facilitando trazabilidad y auditoría \cite{airflow_docs}.

\subsection{Cambio de procesos de arranque}
En Airflow 3, el webserver pasa a ser un API server genérico y el DAG processor se inicia por separado (aun en setups de desarrollo) \cite{airflow_upgrade_3_breaking}. Esto explica por qué ``antes funcionaba con un solo comando'' y ahora hay más procesos: es parte del rediseño arquitectónico.

\section{Mini-caso: envolver un pipeline ETL existente con Airflow (sin redefinir ETL)}
Asumiendo que ya existe un script o librería Python que hace (i) extraer, (ii) transformar y (iii) cargar, el valor de Airflow es \textbf{operar el ciclo completo} con control de errores y auditabilidad. Un patrón típico es:
\begin{enumerate}[leftmargin=*]
  \item \textbf{Preparación}: validar parámetros, verificar conectividad.
  \item \textbf{Ejecución}: correr la carga (por partición, por fecha, por archivo).
  \item \textbf{Calidad}: checks (conteos, nulos, umbrales).
  \item \textbf{Publicación}: marcar un asset/listo, o notificar.
\end{enumerate}

\begin{lstlisting}[style=python,caption={Esqueleto de DAG para orquestar un script ETL existente por fecha}]
from datetime import datetime
from airflow.sdk import dag, task, get_current_context

@dag(
    dag_id="orquesta_etl_por_fecha",
    start_date=datetime(2025, 1, 1),
    schedule="0 2 * * *",
    catchup=False,
    tags=["etl", "orquestacion"]
)
def orquesta():
    @task
    def parametros():
        ctx = get_current_context()
        # En Airflow, la fecha lógica del run sirve como partición de negocio
        logical_date = ctx["logical_date"]  # datetime
        return {"fecha": logical_date.strftime("%Y-%m-%d")}

    @task(retries=2)
    def ejecutar_etl(params: dict):
        fecha = params["fecha"]
        # Llamar librería interna o script (ej: subprocess, import, etc.)
        # run_pipeline(fecha=fecha)
        print(f"Ejecutando ETL para {fecha}")

    @task
    def checks(params: dict):
        fecha = params["fecha"]
        # checks de calidad y auditoría: conteos, duplicados, nulos, etc.
        print(f"Checks de calidad para {fecha}")

    p = parametros()
    ejecutar_etl(p) >> checks(p)

orquesta()
\end{lstlisting}
El ejemplo enfatiza una idea: \textbf{la partición (fecha lógica) vive en el run}, y las tareas son idempotentes y re-ejecutables. Esto habilita reintentos, backfills y diagnósticos de forma natural.

\section{Buenas prácticas y anti-patrones}
\begin{longtable}{p{0.48\textwidth}p{0.48\textwidth}}
\toprule
\textbf{Buena práctica} & \textbf{Anti-patrón típico}\\
\midrule
\endhead
Tareas idempotentes y con efectos controlados. & Tareas que ``acumulan'' datos sin control cuando se reintentan.\\
Usar \texttt{airflow.sdk} y providers; evitar imports internos. & Importar modelos internos o sesiones del metadata DB (bloqueado en Airflow 3).\\
XCom solo para metadatos (rutas, IDs, banderas). & Enviar grandes datasets por XCom.\\
Validar dependencias y credenciales como paso inicial. & Descubrir que falta una librería/permiso en el último paso del pipeline.\\
Minimizar trabajo en parseo del DAG. & Consultas o llamadas API a nivel global del archivo DAG.\\
Usar TaskGroup / mapping para legibilidad y escalabilidad. & DAGs monolíticos con decenas de tareas sin estructura.\\
\bottomrule
\end{longtable}

\section{Conclusiones}
Airflow 3.x no es ``otra herramienta ETL'', sino una plataforma para \textbf{operar workflows} con disciplina: planificación, dependencias, reintentos, trazabilidad y control de ejecuciones históricas. La versión 3 introduce cambios relevantes (API server, DAG processor independiente, Task Execution API y \texttt{airflow.sdk} como interfaz estable) que influyen directamente en cómo se escriben DAGs, cómo se accede a configuración (Variables/Connections) y cómo se mantiene seguridad del sistema \cite{airflow_release_notes_3,airflow_upgrading_3_stable,airflow_public_interface}.

Como siguiente paso práctico (en laboratorio), la recomendación es tomar un pipeline ya construido en Python, envolverlo en un DAG con particionamiento por fecha, agregar checks de calidad y practicar backfills controlados. Esta transición del ``script que corre'' al ``proceso que se opera'' es el objetivo central de aprender Airflow.

\newpage
\printbibliography
\end{document}
